\section{\sys{}}
\label{sec:system}

\begin{table}[t]
\centering\scriptsize
\renewcommand{\arraystretch}{1.05}
\begin{tabular}{@{}p{17mm}p{54.5mm}@{}}
\toprule
\textbf{Term} & \textbf{Definition} \\
\midrule
LLM, model & A large language model: a neural network that writes text. \\
Agent & A program in which a model selects the next action, for example a tool call, until the
task ends. \\
Tool & A function that an agent can call, for example a search or a database query. \\
Workflow & A program that connects model calls and tool calls to do a task. \\
Task & One input to a workflow, with a known correct result. \\
Evaluation & A run of a workflow on a set of tasks. It gives a score for each task. \\
Prompt & The text that a model receives as its instruction. \\
Runtime & The part of \sys{} that receives each call of a workflow and applies the recipe. \\
Trace & The record of all calls and outputs of one run. \\
Token & A small unit of text. Providers charge for each token. \\
API & The interface through which a program calls a model provider. \\
Variant & One version of a workflow. An \emph{edit} tries to improve a prompt. A \emph{mutant} is a
variant that we break on purpose. \\
Miniature & A copy of a workflow with shortcuts on its calls. \\
Shortcut, recipe & A shortcut changes how the runtime handles a call. A recipe is the set of
shortcuts of one miniature. \\
Replay & To use the recorded output of a call, not a new call. \\
Simulated tool & A model writes the output of a tool. A \emph{pure} simulator invents each
output. A \emph{grounded} simulator uses recorded real outputs first. \\
Merged call & One model call writes a full agent run and imagines each tool result. \\
Stub & One model call that replaces a full sub-agent. \\
W3 roles & The planner splits the question. Researchers search Wikipedia. The writer answers. The reviewer finds errors, and the reviser corrects them. \\
Judge & A separate model that grades each run and gives a reason for each failure. \\
Standard error & The expected size of the random error in a mean score. \\
Clear decision & A pair of variants whose full scores differ by more than two standard errors. \\
Noise ceiling & The agreement of the full workflow with a second full run of itself. \\
Kendall's $\tau$ & Agreement of two rankings, from $-1$ (opposite order) to $1$ (same order). \\
Pearson's $r$ & Linear agreement of two lists of numbers, from $-1$ to $1$. \\
Scott's $\pi$ & Agreement of two lists of labels. $0$ is chance and $1$ is full agreement. \\
\bottomrule
\end{tabular}
\caption{Technical terms that are not in the STE dictionary.}
\label{tab:terms}
\end{table}


Table~\ref{tab:4w1h} describes \sys{} in the 4W1H form: what, why, when, where and how. This
section gives more detail on its three parts: the interception point, the shortcuts and the
validity check.

\paragraph{One interception point.}
\sys{} does not change the code of a workflow. Each model call and each tool call goes through one
runtime, and a recipe tells the runtime what to do with each call. Thus, one workflow can give many
miniatures, and each miniature runs the same code as the full workflow.

\paragraph{Honest simulation.}
Two design rules keep the simulated parts honest. A simulator sees only the description of the tool
and the arguments of the call. It never sees the task, because a real tool does not see the task.
Also, the runtime keeps each simulated output, so all variants see the same simulated world.

\begin{table}[t]
\centering\scriptsize
\renewcommand{\arraystretch}{1.05}
\begin{tabular}{@{}l@{\hspace{5pt}}p{19mm}p{47.5mm}@{}}
\toprule
& \textbf{Shortcut} & \textbf{What it does} \\
\midrule
\multicolumn{3}{@{}l}{\textcolor{redblue}{\textbf{A. Replay}}: no loss when the inputs are the same} \\
\op{A1} & Prefix replay & Steps before the changed part get the same inputs as in a recorded run.
Use the recorded outputs. \\
\op{A2} & Tool cache & Use the recorded output of the same tool call. \\
\op{A3} & Paired runs & Run each variant on the same tasks, and compare task by task. \\
\addlinespace
\multicolumn{3}{@{}l}{\textcolor{redblue}{\textbf{B. Reduce}}: less work, low risk} \\
\op{B1} & Slice & Do not run parts whose output does not reach the judge. \\
\op{B2} & Step cap & Limit the number of agent steps. \\
\op{B3} & Cheaper model & Use a cheaper model for the parts that you do not change. \\
\op{B4} & Fewer tasks & Use a fixed random subset of the tasks. \\
\addlinespace
\multicolumn{3}{@{}l}{\textcolor{simgreen!80!black}{\textbf{C. Simulate}}: most savings, check first} \\
\op{C1} & Simulated tools & A model plays the tool (pure or grounded). \\
\op{C2} & Merged calls & One call writes the full agent run, with imagined tool results. \\
\op{C3} & Shorter prompts & Compress the fixed part of a prompt. \\
\op{C4} & Stubbed sub-agent & One call replaces a full sub-agent. \\
\op{C5} & Early stop & Stop the run after the part that you test. \\
\op{C6} & Short outputs & Cut each tool output to a fixed length. \\
\bottomrule
\end{tabular}
\caption{The thirteen shortcuts. A recipe can combine shortcuts from all groups.}
\label{tab:shortcuts}
\end{table}


\paragraph{Thirteen shortcuts.}
Table~\ref{tab:shortcuts} lists the shortcuts in three groups. Group~A replays recorded outputs: if the
inputs of a call do not change, its output does not change. Thus, group~A loses no information. Group~B does less work: fewer tasks, fewer steps, or a cheaper model
for the parts that you do not test. Group~C lets a model play a part of the workflow. It saves the
most work, but it can also change the behaviour that you want to test.

\paragraph{The validity check.}
A miniature is valid for a team if it makes the same decisions as the full workflow. We measure
this in five ways:
\begin{itemize}
  \item \textbf{Clear decisions kept}, our main measure. A pair of variants is \emph{clear} if
  their full scores differ by more than two standard errors. We count the clear pairs where the
  miniature picks the same winner.
  \item \textbf{Ranking agreement}: Kendall's $\tau$ between the two rankings of all variants.
  \item \textbf{Mutants caught}: the miniature must also find the variants that we broke.
  \item \textbf{Reason agreement}: Scott's $\pi$ between the failure reasons of the judge, run by
  run.
  \item \textbf{Fix transfer}: fixes written from the failures of the miniature must also improve
  the full workflow.
\end{itemize}

\paragraph{The noise ceiling.}
We read each measure against the noise ceiling: the same measure between the full workflow and a
second full run. A workflow is random, and test sets are small. Thus, the full workflow does not
fully agree with itself. A miniature cannot do better than this ceiling, except by chance.
